\section{Exercitii}

Facem cativa algoritmi de baza, din cei pe care ii stii deja din C++, dar
scrisi acum in Python, de mana, cu bucla explicita. Daca ai mai citit prin
preajma, stii ca nu iti trebuie neaparat forma asta, o sa vedem in
\hyperref[chapter:numpy-pandas]{capitolul urmator} ca NumPy face acelasi
lucru intr-o singura linie. Dar e important sa trecem prin ele de mana,
macar o data, inainte sa le facem cu o librarie: altfel folosesti
\pycode{numpy.sum()} fara sa stii ce calculeaza de fapt.

\subsection{Suma cifrelor unui numar}
\label{problem:suma-cifrelor}

\href{notebooks/suma-cifrelor.ipynb}{notebook}
$\bullet$
\hyperref[code:suma-cifrelor]{surse}

Extragem cifrele unui numar una cate una, cu impartire si rest la 10, si
le adunam. Genul de exercitiu pe care l-ai facut deja de zeci de ori in
C++, doar sintaxa difera.

\subsection{Suma elementelor unui vector}
\label{problem:suma-vector}

\href{notebooks/suma-vector.ipynb}{notebook}
$\bullet$
\hyperref[code:suma-vector]{surse}

Parcurgem lista element cu element si acumulam suma. O sa refacem exact
exercitiul asta in capitolul urmator, cu \pycode{numpy.sum()}, ca sa vezi
diferenta de cod, si de viteza.

\subsection{Maximul dintr-un vector}
\label{problem:maxim-vector}

\href{notebooks/maxim-vector.ipynb}{notebook}
$\bullet$
\hyperref[code:maxim-vector]{surse}

Aceeasi idee ca la maximul dintr-un vector in C++: retinem cel mai mare
element vazut pana acum, pe masura ce parcurgem lista.

\subsection{Numararea elementelor pare}
\label{problem:numarare-conditie}

\href{notebooks/numarare-conditie.ipynb}{notebook}
$\bullet$
\hyperref[code:numarare-conditie]{surse}

Parcurgem lista si numaram cate elemente indeplinesc o conditie, aici,
sunt pare. E un tipar care revine des: parcurgere plus conditie plus
contor.
